%% A realistic Elsevier CAS manuscript, shaped like the front matter in
%% Elsevier's elsdoc-cas.tex. Used to test venue-shift; not a real paper.
\documentclass[a4paper,fleqn]{cas-dc}

\usepackage[authoryear]{natbib}
\usepackage{pgfplots}

\begin{document}

\shorttitle{Trust Calibration in Automated Driving}
\shortauthors{M. Colley and J. Doe}

\title [mode = title]{Trust Calibration in Highly Automated Driving}
\tnotemark[1]
\tnotetext[1]{This work was funded by the German Research Foundation.}

\author[1]{Mark Colley}[orcid=0000-0001-5207-5029]
\cormark[1]
\ead{mark.colley@uni-ulm.de}
\credit{Conceptualization, Methodology, Writing - original draft}

\author[2]{Jane Doe}
\ead{jane.doe@example.edu}
\ead[url]{https://example.edu/~jdoe}
\credit{Formal analysis, Writing - review \& editing}

\affiliation[1]{organization={Institute of Media Informatics, Ulm University},
                addressline={James-Franck-Ring},
                city={Ulm},
                postcode={89081},
                country={Germany}}

\affiliation[2]{organization={Example University},
                city={Boston},
                state={Massachusetts},
                country={USA}}

\cortext[1]{Corresponding author.}

\begin{abstract}
Automated vehicles must communicate their intent to build appropriate trust.
We report a within-subjects study with 24 participants examining how takeover
requests shape trust over repeated exposure.
\end{abstract}

\begin{highlights}
\item Trust in automated vehicles develops across repeated takeovers.
\item Takeover requests shape trust calibration.
\item A within-subjects study with 24 participants.
\end{highlights}

\begin{keywords}
automated driving \sep trust \sep takeover request \sep user study
\end{keywords}

\maketitle

\section{Introduction}
Trust in automated vehicles is neither uniformly good nor uniformly
bad~\citep{lee2004trust}. We examine how it develops across repeated takeovers.

\begin{figure}[t]
  \centering
  \includegraphics[width=\columnwidth]{figures/apparatus}
  \caption{The driving simulator used in the study.}
  \label{fig:apparatus}
\end{figure}

\section{Method}
We ran a within-subjects experiment with 24 participants.

\begin{table}[t]
  \caption{Participant demographics.}
  \label{tab:demographics}
  \begin{tabular}{lr}
    \toprule
    Measure & Value \\
    \midrule
    Participants & 24 \\
    Mean age & 27.4 \\
    \bottomrule
  \end{tabular}
\end{table}

\section{Results}
Trust increased over exposure, $F(2,46) = 8.1$, $p < .001$.

\section{Conclusion}
Repeated exposure calibrates trust.

\printcredits

\bibliographystyle{cas-model2-names}
\bibliography{references}

\end{document}
